% Options for packages loaded elsewhere
\PassOptionsToPackage{unicode}{hyperref}
\PassOptionsToPackage{hyphens}{url}
\documentclass[
]{article}
\usepackage{xcolor}
\usepackage{amsmath,amssymb}
\setcounter{secnumdepth}{-\maxdimen} % remove section numbering
\usepackage{iftex}
\ifPDFTeX
  \usepackage[T1]{fontenc}
  \usepackage[utf8]{inputenc}
  \usepackage{textcomp} % provide euro and other symbols
\else % if luatex or xetex
  \usepackage{unicode-math} % this also loads fontspec
  \defaultfontfeatures{Scale=MatchLowercase}
  \defaultfontfeatures[\rmfamily]{Ligatures=TeX,Scale=1}
\fi
\usepackage{lmodern}
\ifPDFTeX\else
  % xetex/luatex font selection
\fi
% Use upquote if available, for straight quotes in verbatim environments
\IfFileExists{upquote.sty}{\usepackage{upquote}}{}
\IfFileExists{microtype.sty}{% use microtype if available
  \usepackage[]{microtype}
  \UseMicrotypeSet[protrusion]{basicmath} % disable protrusion for tt fonts
}{}
\makeatletter
\@ifundefined{KOMAClassName}{% if non-KOMA class
  \IfFileExists{parskip.sty}{%
    \usepackage{parskip}
  }{% else
    \setlength{\parindent}{0pt}
    \setlength{\parskip}{6pt plus 2pt minus 1pt}}
}{% if KOMA class
  \KOMAoptions{parskip=half}}
\makeatother
\setlength{\emergencystretch}{3em} % prevent overfull lines
\providecommand{\tightlist}{%
  \setlength{\itemsep}{0pt}\setlength{\parskip}{0pt}}
\usepackage{bookmark}
\IfFileExists{xurl.sty}{\usepackage{xurl}}{} % add URL line breaks if available
\urlstyle{same}
\hypersetup{
  hidelinks,
  pdfcreator={LaTeX via pandoc}}

\author{}
\date{}

\begin{document}

\textbf{RNR Coding --- Reference Implementation Conformance
Specification}

\emph{Engineering companion to the framework paper, v0.8}

Pavlo Ratushnyi --- May 2026

\textbf{1. Purpose and scope}

This document specifies the conformance requirements for a reference
implementation of RNR Coding. It operationalizes Theorems 10.1, 10.2,
10.3, 10.5, and 10.6 of the main paper into a software-engineering
specification that can be implemented, tested, and certified across the
hardware targets named in §10.6.

The specification is RFC-style: requirements use the MUST/SHOULD/MAY
language of RFC 2119. A conforming implementation MUST satisfy every
MUST-level requirement on every selected hardware target. SHOULD-level
requirements may be omitted with documented rationale; MAY-level
requirements are at implementer discretion.

\textbf{1.1 In scope}

\begin{itemize}
\item
  The integer-arithmetic specification under which Theorem
  10.1\textquotesingle s bit-exact reproduction holds.
\item
  Non-linearity lookup table format and serialization.
\item
  Per-hardware-target conformance specifications (CPU scalar, CPU SIMD
  including AVX-512 and NEON, GPU integer paths, NPU paths including
  Apple Neural Engine, Google Edge TPU, Qualcomm Hexagon, AWS
  Inferentia).
\item
  Sync-point and seek-index binary formats supporting Theorem 10.3
  random access.
\item
  Block-device driver architecture supporting Theorem 10.5
  mounted-filesystem deployment.
\item
  Conformance test suite with explicit acceptance criteria.
\end{itemize}

\textbf{1.2 Out of scope}

\begin{itemize}
\item
  Predictor training. This specification accepts a trained predictor as
  input; how the predictor is trained is upstream.
\item
  Specific predictor architectures. The specification is
  architecture-neutral within the constraints of §3.
\item
  Storage encryption and content addressing. These are deployment-layer
  concerns and not part of the codec.
\item
  Network distribution protocols and authentication. Deployment-layer.
\end{itemize}

\textbf{1.3 Relation to the main paper}

The main paper proves the mathematical results; this document specifies
how to instantiate them in code. No new theorems are introduced. Every
requirement in this document is traceable to a theorem in the main
paper; the traceability matrix is given in §14.

\textbf{2. Terminology}

\textbf{MUST, REQUIRED.} An absolute requirement. A conforming
implementation MUST satisfy this.

\textbf{SHOULD, RECOMMENDED.} Valid implementations may omit this with
documented rationale, but the recommended path produces better
reliability or interoperability.

\textbf{MAY, OPTIONAL.} Implementer discretion. The implementation is
conforming whether or not the option is taken.

\textbf{Reference target.} The CPU scalar implementation (§4.1) is the
canonical reference. All other targets MUST produce output bit-identical
to the reference for any input.

\textbf{Conforming target.} A hardware target whose implementation
satisfies all MUST-level requirements applicable to that target class
and passes the conformance test suite of §12.

\textbf{3. Predictor inference engine: universal requirements}

Per Theorem 10.1 of the main paper, every conforming predictor inference
engine MUST satisfy the following requirements. These are universal
across hardware targets and apply before any target-specific
specialization.

\textbf{R-3.1 (Integer arithmetic).} All intermediate computations in
the inference path MUST use integer arithmetic with bit-widths specified
by the model description. No floating-point arithmetic is permitted in
the inference path except as allowed in R-3.6.

\textbf{R-3.2 (Overflow avoidance).} The accumulator bit-width MUST
be at least
$\lceil \log_2(L \cdot W_{\max} \cdot I_{\max}) \rceil + 1$, where
$L$ is the maximum layer depth, $W_{\max}$ is an upper bound on the
$L_1$ norm of any weight row, and $I_{\max}$ is an upper bound on
the magnitude of any layer input. The bounds $W_{\max}$ and
$I_{\max}$ MUST be declared in the model description; an
implementation MAY verify them at model-load time.

\textbf{R-3.3 (Non-linearities).} All non-linearity operations MUST be
implemented as deterministic integer-to-integer lookup tables specified
in the model description per §6. Approximations computed at runtime
(e.g., Taylor series, polynomial fits) are forbidden.

\textbf{R-3.4 (Reduction order).} All reduction operations (sums in
matrix-vector products, softmax denominators, normalization aggregates)
MUST follow the order specified by the model description. The order is
encoded as a tree shape; implementations MUST honor that tree shape
exactly. Platform-default reduction orders (GPU atomic reductions,
randomized parallel reductions, OpenMP reduce with default schedule) are
forbidden.

\textbf{R-3.5 (Determinism).} All operations in the inference path MUST
be deterministic. No randomness is permitted, including random
initialization, dropout, or stochastic rounding. If the predictor
architecture requires randomness for training, the training-mode
randomness MUST be disabled at inference time.

\textbf{R-3.6 (Initialization).} Floating-point arithmetic MAY be used
during model description parsing (loading weights from a model file) but
MUST produce only integer outputs that are subsequently used by
R-3.1-conforming operations. No floating-point value computed at parse
time may flow into the inference path.

\textbf{R-3.7 (Layer ordering).} Layer execution order MUST follow the
topological order specified by the model description. Implementations
MAY parallelize within a layer but MUST NOT reorder layer-level
dependencies.

\textbf{R-3.8 (No fused-multiply-add unless specified).} FMA
instructions (FMA3, FMA4) on floating-point operands are forbidden.
Integer FMA-equivalents (VPDPWSSD, VPDPBUSD, MFMA-int, Tensor Core
int8/int4 paths) are allowed where the integer arithmetic semantics are
bit-equivalent to scalar integer multiply-add.

\textbf{4. Hardware-target specifications}

A conforming implementation MUST target at least one of the hardware
paths below. Each path MUST satisfy R-3.1 through R-3.8 in addition to
its target-specific requirements. Cross-target byte-identity is required
by Theorem 10.6 and verified by the conformance suite of §12.

\textbf{4.1 CPU scalar reference target}

The CPU scalar target is the canonical reference. All other targets MUST
produce bit-identical output to this target on the conformance test
vectors of §12.

\textbf{R-4.1.1 (Default types).} Weights are int8 (signed, range -128
to 127). Activations are int32 (signed). Accumulators are int64
(signed). The model description MAY override per-layer.

\textbf{R-4.1.2 (Multiplication).} Weight-activation multiplication is
int32 * int8 -\textgreater{} int64 with explicit sign extension. The C99
expression is: (int64\_t)(int32\_t)activation *
(int64\_t)(int8\_t)weight.

\textbf{R-4.1.3 (Accumulation).} Sequential left-to-right accumulation
in int64. The compiler MUST NOT reorder additions; implementations MUST
use explicit accumulator variables, not standard library reductions
whose ordering is implementation-defined.

\textbf{R-4.1.4 (Output scaling).} After accumulation, the int64
accumulator is shifted right by the layer\textquotesingle s output-scale
parameter (specified in the model description) and clamped to the
activation type range. The shift is arithmetic (sign-extending).

\textbf{R-4.1.5 (Lookup tables).} Non-linearity tables are stored as
flat int32 arrays. The index is the (clamped) input value; the value is
the (signed) output. Bounds checking MUST be present in debug builds and
SHOULD be present in release builds (modern branch predictors make the
check nearly free).

\textbf{4.2 CPU SIMD targets}

\emph{\textbf{4.2.1 AVX2}}

\textbf{R-4.2.1.1 (Vector width).} 256-bit vectors holding 8 lanes of
int32 or 32 lanes of int8.

\textbf{R-4.2.1.2 (Multiply-accumulate).} MUST use \_mm256\_madd\_epi16
for int16*int16-\textgreater int32 fused multiply-add. If AVX-VNNI is
available, MAY use \_mm256\_dpwssd\_avx\_epi32 for int8 dot products.

\textbf{R-4.2.1.3 (Horizontal reduction).} MUST be implemented as a
sequence of \_mm256\_hadd\_epi32 (or shuffles plus adds) following the
tree shape specified by the model description. Implementations MUST NOT
use \_mm256\_reduce\_add\_epi32 unless the
implementation\textquotesingle s reduction tree is documented and
matches the model description.

\emph{\textbf{4.2.2 AVX-512}}

\textbf{R-4.2.2.1 (Vector width).} 512-bit vectors holding 16 lanes of
int32 or 64 lanes of int8.

\textbf{R-4.2.2.2 (VNNI instructions).} MUST use VPDPBUSD for
int8*int8-\textgreater int32 and VPDPWSSD for
int16*int16-\textgreater int32 when these instructions are available
(Cascade Lake and later for VNNI; Ice Lake and later for IFMA).

\textbf{R-4.2.2.3 (Mask registers).} Predication via opmask registers
(k1-k7) is permitted. Operations using opmask MUST produce the same
output as scalar operations on the corresponding lanes.

\emph{\textbf{4.2.3 ARM NEON}}

\textbf{R-4.2.3.1 (Vector width).} 128-bit vectors holding 4 lanes of
int32 or 16 lanes of int8.

\textbf{R-4.2.3.2 (Multiply-accumulate).} MUST use VMLAL (int16
multiply, int32 accumulate) or SDOT (int8 dot product when
ARMv8.2-DotProd is available).

\textbf{R-4.2.3.3 (Horizontal reduction).} MUST use VADDV with
sequential ordering or explicit tree reductions matching the model
description.

\emph{\textbf{4.2.4 RISC-V Vector}}

\textbf{R-4.2.4.1 (Vector configuration).} VL and LMUL MUST be
configured deterministically per layer; the configuration MUST be
specified in the model description. Auto-vectorization that selects VL
based on hardware preferences is forbidden in the conforming path.

\textbf{R-4.2.4.2 (Integer instructions).} MUST use vmul.vv / vredsum.vs
for multiply-reduce sequences. The reduction follows
RVV\textquotesingle s specified left-fold ordering.

\textbf{4.3 GPU integer paths}

\emph{\textbf{4.3.1 NVIDIA Tensor Core path}}

\textbf{R-4.3.1.1 (Deterministic mode).} cuBLASLt MUST be in
deterministic mode (CUBLASLT\_DETERMINISTIC\_MODE = 1). cuDNN MUST be
configured with the deterministic algorithm hint
(cudnnSetDeterministic). Algorithm selection MUST be pinned in the model
description.

\textbf{R-4.3.1.2 (Tensor Core paths).} Integer multiply-accumulate MUST
use Tensor Core int8 (HMMA.16816.S8) or int4 (HMMA.16832.S4) MMA
instructions on Turing or later. Mixed-precision (int8 * fp16) Tensor
Core paths are forbidden.

\textbf{R-4.3.1.3 (Reduction tree).} Tensor Core MMA produces partial
sums per warp; the cross-warp reduction MUST follow the tree specified
by the model description, implemented via shared memory and
\_\_syncwarp(). Atomic reductions are forbidden.

\textbf{R-4.3.1.4 (Kernel scheduling).} The kernel launch configuration
(grid size, block size, shared memory) MUST be derived deterministically
from the model description. Library auto-tuners MUST be disabled.

\emph{\textbf{4.3.2 AMD CDNA matrix path}}

\textbf{R-4.3.2.1 (MFMA instructions).} Integer matrix
multiply-accumulate MUST use V\_MFMA\_I32\_xxx instructions.
Mixed-precision MFMA paths are forbidden.

\textbf{R-4.3.2.2 (rocBLAS configuration).} rocBLAS MUST be in
deterministic mode. Algorithm selection MUST be pinned in the model
description.

\textbf{R-4.3.2.3 (Cross-CU reduction).} Reductions across compute units
MUST follow the model description\textquotesingle s tree. Atomic
reductions are forbidden.

\textbf{4.4 NPU paths}

\emph{\textbf{4.4.1 Apple Neural Engine (ANE)}}

\textbf{R-4.4.1.1 (Core ML integer path).} MUST use the int8 inference
subset of Core ML 8.0 or later, with quantized model format (mlpackage
with quantized weights). Float16 paths are forbidden.

\textbf{R-4.4.1.2 (Reduction order).} ANE\textquotesingle s reduction
order is vendor-fixed. Conformance requires verifying that
ANE\textquotesingle s reduction matches the CPU scalar reference for the
specific model architecture; if it does not, the model description MUST
specify an alternative reduction order achievable on ANE, or ANE MUST be
excluded from the conforming targets for that model.

\textbf{R-4.4.1.3 (Precision verification).} The conformance test of §12
MUST be run on ANE before certification. Test vectors of at least 10,000
inputs MUST produce byte-identical predictor output to the CPU
reference.

\emph{\textbf{4.4.2 Google Edge TPU}}

\textbf{R-4.4.2.1 (TensorFlow Lite micro path).} MUST use the TensorFlow
Lite quantized model format with int8 weights and int8 activations.
Per-tensor or per-channel quantization MUST match the model description.

\textbf{R-4.4.2.2 (Reduction).} Edge TPU\textquotesingle s reduction
order is fixed and known to be deterministic across firmware versions
documented in the model description\textquotesingle s compatibility
manifest.

\emph{\textbf{4.4.3 Qualcomm Hexagon NPU}}

\textbf{R-4.4.3.1 (HVX integer path).} MUST use Hexagon Vector eXtension
(HVX) integer instructions with explicit accumulator typing. The Hexagon
NN library MUST be configured in deterministic mode.

\textbf{R-4.4.3.2 (HMX support).} If the target supports Hexagon Matrix
eXtension (HMX), MUST use HMX int8 path with pinned tile sizes.

\emph{\textbf{4.4.4 AWS Inferentia}}

\textbf{R-4.4.4.1 (Neuron SDK).} MUST use the AWS Neuron SDK 2.0 or
later, with int8 quantization. The Neuron compiler MUST be invoked with
deterministic flags.

\textbf{R-4.4.4.2 (Tensor partitioning).} Tensor partitioning across
NeuronCores MUST be pinned in the model description. Auto-partitioning
is forbidden in the conforming path.

\emph{\textbf{4.4.5 Common NPU requirement}}

\textbf{R-4.4.5 (Cross-target verification).} Each NPU target MUST pass
the conformance test suite of §12 on a fixed 10,000-input test vector
before being certified as conforming. The test vector MUST include all
non-linearity types used by the predictor architecture.

\textbf{5. Integer-arithmetic specification}

\textbf{5.1 Bit-widths}

\textbf{R-5.1.1 (Defaults).} Default weight bit-width is int8 (signed,
range {[}-128, 127{]}). Default activation bit-width is int32. Default
accumulator bit-width is int64. Default lookup-table value bit-width is
int32.

\textbf{R-5.1.2 (Model overrides).} The model description MAY override
defaults per-layer. Common overrides: int4 weights for high compression
at slightly reduced accuracy; int16 activations for higher dynamic
range.

\textbf{5.2 Overflow behavior}

\textbf{R-5.2.1 (Defined behavior).} Implementations MUST NOT rely on
undefined-behavior wraparound. All signed integer arithmetic MUST be in
well-defined two\textquotesingle s-complement representation; on
platforms where the C compiler may interpret signed overflow as
undefined behavior, the implementation MUST either use unsigned
arithmetic with explicit sign conversion, or compile with -fwrapv (gcc)
or equivalent.

\textbf{R-5.2.2 (No saturation by default).} Arithmetic operations MUST
NOT saturate unless explicitly specified per-layer in the model
description. Saturation is a layer property, not a hardware property.

\textbf{5.3 Division and shift}

\textbf{R-5.3.1 (Division).} Signed integer division MUST follow C99
truncation-toward-zero semantics. Floor division (Python-style) is
forbidden in the conforming path.

\textbf{R-5.3.2 (Modulo).} Signed integer modulo MUST follow C99 (sign
of the result matches the dividend). Mathematical modulo (always
non-negative) is forbidden.

\textbf{R-5.3.3 (Shift).} Arithmetic shift right MUST be sign-extending.
The shift amount MUST be in {[}0, bit\_width - 1{]}; out-of-range shifts
are forbidden and MUST be caught at model-load time.

\textbf{5.4 Reduction order}

\textbf{R-5.4.1 (Specification).} The model description MUST specify
reduction order as one of: (a) sequential-left, (b) sequential-right, or
(c) tree with explicit tree shape (an array of pair-merge indices).

\textbf{R-5.4.2 (Implementation).} Implementations MUST produce
identical results to the specified reduction order. Hardware that cannot
directly produce the order MUST emulate it (e.g., by sequential
left-fold on a single SIMD lane).

\textbf{6. Non-linearity lookup table specification}

All non-linearity operations are integer-to-integer lookup tables
specified in the model description. The lookup table format is canonical
across hardware targets.

\textbf{6.1 Table format}

\textbf{R-6.1.1 (Layout).} Each table is a contiguous array of int32
values. The array is indexed by the (clamped) input value, offset to be
non-negative.

\textbf{R-6.1.2 (Header).} Each table is preceded by a 16-byte header:
4-byte magic \textquotesingle RNRT\textquotesingle, 4-byte version,
4-byte input range minimum, 4-byte input range maximum. The array length
is (max - min + 1).

\textbf{R-6.1.3 (Clamping).} Inputs outside {[}min, max{]} MUST be
clamped to the range before lookup. Clamping is symmetric: values below
min clamp to min, values above max clamp to max.

\textbf{6.2 Common non-linearities}

The following non-linearities have canonical specifications. Conforming
implementations MUST use these specifications when the corresponding
non-linearity is named in the model description.

\emph{\textbf{6.2.1 Sigmoid}}

$\mathrm{Sigmoid}(x) = 1 / (1 + \exp(-x))$. The integer table is
parameterized by an input scale $s_{\text{in}}$ and output scale
$s_{\text{out}}$:
$\mathrm{table}[i] = \mathrm{round}(\mathrm{sigmoid}(i / s_{\text{in}}) \cdot s_{\text{out}})$.
The model description specifies $s_{\text{in}}$ (default 256) and
$s_{\text{out}}$ (default 65536).

\emph{\textbf{6.2.2 Tanh}}

$\tanh(x) = (\exp(x) - \exp(-x)) / (\exp(x) + \exp(-x))$. Similar
parameterization to sigmoid; default $s_{\text{in}} = 128$,
$s_{\text{out}} = 32768$.

\emph{\textbf{6.2.3 GELU}}

$\mathrm{GELU}(x) = x \cdot \Phi(x)$, where $\Phi$ is the standard
normal CDF. Parameterized by $s_{\text{in}}$ (default 256) and
$s_{\text{out}}$ (default 65536).

\emph{\textbf{6.2.4 ReLU}}

$\mathrm{ReLU}(x) = \max(0, x)$. MUST be implemented as
$\max(0, x)$ directly, not as a table; this is a special case for
efficiency.

\emph{\textbf{6.2.5 Softmax}}

Softmax requires special handling because it involves division and
is not directly a table lookup. The decomposition is:
\begin{enumerate}
\item shift: $x' = x - \max(x)$ \quad (numerical stability)
\item exp\_lookup: $y_i = \mathrm{exp\_table}[x'_i]$ \quad (integer table)
\item sum: $S = \sum_i y_i$ \quad (ordered reduction)
\item division: $z_i = (y_i \cdot \mathrm{scale}) / S$ \quad (integer division)
\end{enumerate}
Steps 2, 3, and 4 each follow R-3.3, R-3.4, and R-5.3.1 respectively.

\emph{\textbf{6.2.6 Layer normalization}}

Layer normalization requires integer reciprocal-square-root for the
variance scaling. The decomposition is:
\begin{enumerate}
\item mean: $\mu = \sum_i x_i / n$ \quad (ordered reduction + division)
\item centered: $c_i = x_i - \mu$
\item variance: $v = \sum_i c_i^2 / n$ \quad (ordered reduction + division)
\item rsqrt\_lookup: $r = \mathrm{rsqrt\_table}[v + \varepsilon]$
\item scale: $z_i = c_i \cdot r \cdot \gamma + \beta$
\end{enumerate}
The rsqrt table covers the variance range observed in training, with
$\varepsilon$ added to avoid division by zero. $\varepsilon$ and the
rsqrt table parameters MUST be in the model description.

\textbf{7. Adaptive predictor conformance (Theorem 10.2)}

Implementations supporting online adaptation MUST use a compatible
online optimizer per Definition 10.1 of the main paper. The following
optimizers are pre-certified compatible.

\textbf{7.1 Integer SGD}

\textbf{R-7.1.1 (Update rule).} The update is
$\text{param}[i] \mathrel{-}= \mathrm{shift\_right}(\text{grad}[i] \cdot \mathrm{learning\_rate}, \mathrm{lr\_scale})$.

\textbf{R-7.1.2 (Parameter ranges).} $\mathrm{learning\_rate}$ MUST
be a non-negative integer. $\mathrm{lr\_scale}$ MUST be a
non-negative integer specifying the right-shift after multiplication.

\textbf{R-7.1.3 (Saturation).} If the update produces a value
outside the parameter type range, the parameter MUST be clamped to
the type range. The clamping behavior MUST match the CPU reference
target.

\textbf{7.2 SignSGD}

\textbf{R-7.2.1 (Update rule).} The update is
$\text{param}[i] \mathrel{-}= \mathrm{learning\_rate} \cdot \mathrm{sign}(\text{grad}[i])$,
where $\mathrm{sign}$ returns $-1, 0$, or $+1$ for negative, zero,
and positive gradient respectively.

\textbf{R-7.2.2 (Trivial conformance).} SignSGD is trivially
compatible: the update function is integer-to-integer with no
real-valued operations.

\textbf{7.3 Integer Adam}

\textbf{R-7.3.1 (State).} Adam state per parameter is $(m, v)$ where
$m$ is the first moment and $v$ is the second moment, both
integer-valued.

\textbf{R-7.3.2 (Update rule).}
$m_{\text{new}} = \beta_1 \cdot m + (1 - \beta_1) \cdot \text{grad}$
(in integer arithmetic with explicit scaling);
$v_{\text{new}} = \beta_2 \cdot v + (1 - \beta_2) \cdot \text{grad}^2$;
$\text{param}_{\text{new}} = \text{param} - \mathrm{learning\_rate} \cdot m_{\text{new}} \cdot \mathrm{rsqrt\_table}[v_{\text{new}} + \varepsilon]$.
The $\mathrm{rsqrt\_table}$ is part of the model description.

\textbf{R-7.3.3 (Reciprocal-square-root table).} The integer
rsqrt\_table for Adam MUST cover the second-moment range observed in
training. The table format follows §6.1.

\textbf{7.4 Gradient computation}

\textbf{R-7.4.1 (Integer backpropagation).} Backpropagation through the
inference graph MUST satisfy R-3.1 through R-3.8 (the same
integer-arithmetic constraints as forward inference).

\textbf{R-7.4.2 (Loss function).} The loss function used for adaptation
MUST be specified in the model description and MUST be implementable as
an integer-to-integer operation given integer model outputs.

\textbf{7.5 Adaptation schedule}

\textbf{R-7.5.1 (Subblock order).} The order in which verified subblocks
are used for adaptation MUST be specified in the model description.
Data-dependent ordering is forbidden.

\textbf{R-7.5.2 (Adaptation frequency).} The model description MUST
specify how often adaptation occurs (e.g., every K-th subblock). Encoder
and decoder MUST adapt at the same positions.

\textbf{8. Random-access conformance (Theorem 10.3)}

\textbf{8.1 Sync-point format}

\textbf{R-8.1.1 (Per-sync-point record).} Each sync point is a record:
(byte\_offset\_uint64, repair\_bit\_offset\_uint64,
snapshot\_length\_uint32, snapshot\_bytes\_variable). For W-bounded
predictors, snapshot\_length is 0 and snapshot\_bytes is empty.

\textbf{R-8.1.2 (Snapshot serialization).} For non-W-bounded predictors,
the snapshot is the predictor\textquotesingle s full internal state,
serialized in the order specified by the model description. The
serialization MUST be canonical: identical state produces identical
bytes on any conforming implementation.

\textbf{8.2 Seek-index format}

\textbf{R-8.2.1 (Index structure).} The seek index is a sorted array of
sync-point records, sorted by byte\_offset ascending. The array length
is encoded as a uint64 prefix to the array.

\textbf{R-8.2.2 (Lookup).} Sync-point lookup for byte position $N$
MUST use binary search to find the largest
$\mathrm{byte\_offset} \leq N$. The lookup is
$O(\log(\mathrm{index\_length}))$.

\textbf{R-8.2.3 (Compression).} The seek index MAY be compressed with a
separate dictionary-based compressor (e.g., zstd) for storage
efficiency. Decompressed form MUST match the canonical layout.

\textbf{8.3 Range-decode function}

The standardized range-decode API:
\begin{verbatim}
function range_decode(archive, N, k):
    # archive: opened RNR archive handle
    # N: byte offset in source
    # k: number of bytes to decode
    # returns: bytes [N, N+k) of the source

    P = binary_search(archive.seek_index, N)  # largest sync_point <= N
    state = initialize_predictor_state(archive, P)
    bit_offset = archive.seek_index[P].repair_bit_offset
    output = bytearray(k)

    for i in [P.byte_offset .. N + k):
        Y_i = predictor(state, context_window)
        repair_bits = read_bits(archive.repair_stream, bit_offset,
                                length_for_position_i)
        bit_offset += length_for_position_i
        byte_i = decode_position(Y_i, repair_bits)
        context_window.append(byte_i)
        if i >= N and i < N + k:
            output[i - N] = byte_i

    return output
\end{verbatim}

\textbf{9. Block-device driver specification (Theorem 10.5)}

\textbf{9.1 Driver API}

The minimal API surface for a block-device driver wrapping an RNR
archive:
\begin{verbatim}
function open(archive_path)                       -> handle
function close(handle)
function read(handle, byte_offset, length, buffer) -> bytes_read
function get_size(handle)                          -> archive_logical_size
\end{verbatim}
Optional extensions for filesystem use:
\begin{verbatim}
function readv(handle, iovecs)              -> bytes_read    # scatter-gather read
function readahead(handle, byte_offset, length)              # cache warmup hint
function fstat(handle)                      -> archive_metadata
\end{verbatim}

\textbf{9.2 Implementation pseudocode}
\begin{verbatim}
function read(handle, byte_offset, length, buffer):
    # Check cache first
    cached = cache_lookup(byte_offset, length)
    if cached:
        copy(cached, buffer)
        return length

    # Decode required range
    chunk = range_decode(handle.archive, byte_offset, length)
    cache_insert(byte_offset, chunk)
    copy(chunk, buffer)
    return length
\end{verbatim}

\textbf{9.3 Cache architecture}

\textbf{R-9.3.1 (Cache organization).} The cache SHOULD be organized as
fixed-size pages (e.g., 64 KiB) aligned to sync-point boundaries. This
allows page-granular eviction and reuse.

\textbf{R-9.3.2 (Eviction policy).} Either LRU (simple, predictable) or
ARC (adaptive, better for mixed workloads) is acceptable. The choice
MUST be configurable by the deployment.

\textbf{R-9.3.3 (Concurrency).} The cache MUST be thread-safe. Multiple
concurrent reads MUST not corrupt cache state.

\textbf{9.4 Concurrency}

\textbf{R-9.4.1 (Read parallelism).} Multiple read() calls on the same
handle MUST be safely concurrent. The driver MAY serialize predictor
inference internally (since most inference engines are single-threaded),
but the driver-level API MUST appear parallel to the caller.

\textbf{R-9.4.2 (Inference engine sharing).} If multiple predictor
instances are available (e.g., one per GPU), the driver MAY load-balance
across them. Load balancing MUST be deterministic (e.g., round-robin) to
preserve byte-exact reproducibility.

\textbf{10. Linux FUSE reference architecture}

A reference implementation using FUSE (Filesystem in Userspace) is the
recommended first deployment. FUSE allows the driver to run as an
unprivileged user-space process, which is appropriate for the maturity
level of the implementation.

\textbf{10.1 FUSE operations}

Implement the following FUSE operations:
\begin{verbatim}
getattr(path)                          -> file metadata
readdir(path)                          -> directory entries  # if filesystem image
open(path, flags)                      -> file handle
read(handle, offset, length, buffer)   -> bytes_read
release(handle)                        -> void
statfs(path)                           -> filesystem statistics
\end{verbatim}

\textbf{10.2 Mounting}

Mount syntax for a single-file archive (e.g., disk image):
\begin{verbatim}
rnrfs --archive=/path/to/image.rnr /mnt/target
\end{verbatim}
The mounted target presents the decompressed archive as a single
read-only file:
\begin{verbatim}
/mnt/target/data    # the decompressed source
\end{verbatim}
For ISO images that contain a filesystem (e.g., iso9660), wrap with a
second mount:
\begin{verbatim}
rnrfs --archive=/path/to/ubuntu.iso.rnr /mnt/rnr
mount -t iso9660 -o loop /mnt/rnr/data /mnt/iso
\end{verbatim}

\textbf{10.3 Write-through overlay}

For read-write use, mount the RNR archive read-only and overlay with
overlayfs:
\begin{verbatim}
rnrfs --archive=/path/to/system.rnr /mnt/base
mkdir /mnt/work /mnt/upper
mount -t overlay overlay \
  -o lowerdir=/mnt/base,upperdir=/mnt/upper,workdir=/mnt/work \
  /mnt/merged
\end{verbatim}

Writes go to /mnt/upper (a normal directory on a writable filesystem);
reads consult /mnt/upper first and fall through to /mnt/base on miss.
This is the standard container-image overlay pattern.

\textbf{10.4 Performance tuning}

\begin{itemize}
\item
  Set FUSE max\_read parameter to a multiple of the sync-point spacing
  (e.g., 256 KiB for K = 64 KiB) to allow batched reads.
\item
  Enable writeback caching only if the overlay layer is
  filesystem-backed; the RNR layer itself is read-only.
\item
  Pin the predictor inference engine to a specific NUMA node on
  multi-socket systems to reduce memory latency.
\item
  On Apple Silicon, use the ANE path (via Core ML) rather than CPU; the
  latency difference is approximately 100x.
\end{itemize}

\textbf{11. Verification under random access}

Theorem 9.1 of the main paper covers chained-hash verification under
sequential decoding. For random access, three verification modes are
supported.

\textbf{11.1 Per-sync-point hashes (recommended for archive-style use)}

\textbf{R-11.1.1 (Hash per sync point).} Each sync point's
seek-index entry MAY include a hash of the bytes $[P, P+K)$ of the
source, computed with SHA-256 truncated to 64 bits. The
random-access reader verifies the hash after decoding the region.

\textbf{R-11.1.2 (Overhead).} Adds 8 bytes per sync point. For
$K = 64$ KiB and $N = 1$ GiB, this is $N/K \cdot 8 = 128$ KiB or
about 0.01\% overhead.

\textbf{11.2 Merkle tree (recommended for content-addressable storage)}

\textbf{R-11.2.1 (Tree structure).} Build a binary Merkle tree over the
source bytes with leaves at sync-point granularity. Each internal node
is SHA-256 of its two children. The root is stored in the archive
header.

\textbf{R-11.2.2 (Random access verification).} A reader that
performs random access at position $N$ obtains the Merkle path from
the leaf containing $N$ up to the root (a logarithmic number of
hashes). The reader verifies the path against the stored root.

\textbf{R-11.2.3 (Overhead).} Adds 32 bytes per sync point plus the
tree storage. For $N = 1$ GiB and $K = 64$ KiB, the tree has
approximately 32 KiB sync points and total tree storage is
approximately 2 MiB, or 0.2\% overhead.

\textbf{11.3 Unverified random access (for trusted archives)}

\textbf{R-11.3.1 (Trust boundary).} If the archive is content-addressed
at a higher layer (e.g., the archive itself is named by a hash known to
the reader), per-access verification may be omitted. Implementations
MUST clearly document when this mode is in use; archives without
verification metadata MUST be marked as
\textquotesingle trusted-only\textquotesingle{} in the archive header.

\textbf{12. Conformance test suite}

A conforming implementation MUST pass the conformance test suite
specified below. Implementations that fail any test are not conforming
on the failing target.

\textbf{12.1 Test vectors}

\textbf{R-12.1.1 (Predictor test vector).} The test suite includes a
fixed set of 10,000 input vectors covering the expected input
distribution of the predictor. For each input, the expected output of M
is precomputed by the CPU reference target and stored in a manifest
file.

\textbf{R-12.1.2 (Test sources).} The test suite includes 100 source
blocks of varying size (1 KiB to 16 MiB) and content (text, binary,
image, mixed). Each block has its expected RNR archive precomputed.

\textbf{R-12.1.3 (Test archives).} The test suite includes 10 prebuilt
RNR archives of varying size (10 MiB to 1 GiB) for random-access
testing. Each archive has 1000 (offset, length) query pairs with
expected output bytes.

\textbf{12.2 Acceptance criteria}

\textbf{R-12.2.1 (Predictor byte-identity).} For all 10,000 inputs, the
target\textquotesingle s predictor output MUST be byte-identical to the
reference target\textquotesingle s output.

\textbf{R-12.2.2 (Encode-decode round-trip).} For all 100 source blocks,
encoding with the target and decoding with the target MUST recover the
original block exactly.

\textbf{R-12.2.3 (Cross-target byte-identity).} For all 100 source
blocks, encoding on the target MUST produce a byte-identical archive to
encoding on the reference target.

\textbf{R-12.2.4 (Random access correctness).} For all 1000
random-access queries on the 10 test archives, the returned bytes MUST
match the expected output.

\textbf{R-12.2.5 (Reproducibility under repetition).} Running the same
test 5 times on the same target MUST produce byte-identical results each
run.

\textbf{12.3 CI/CD integration}

A reference CI pipeline runs the conformance suite on every commit:
\begin{verbatim}
build_matrix:
  - target: cpu-scalar         runner: linux-x86_64
  - target: avx512             runner: linux-x86_64-avx512
  - target: neon               runner: linux-arm64
  - target: nvidia-tensor-core runner: linux-x86_64-nvidia
  - target: amd-mfma           runner: linux-x86_64-amd-gpu
  - target: apple-ane          runner: macos-arm64-m2
  - target: edge-tpu           runner: linux-coral-dev-board

for each target in build_matrix:
    run conformance suite on target
    diff target output against reference manifest
    fail if any byte differs
\end{verbatim}

\textbf{13. Versioning and compatibility}

RNR archives carry a format version in the header. Compatibility rules:

\textbf{R-13.1 (Semantic versioning).} The archive format follows
semantic versioning. MAJOR version changes break archive compatibility;
MINOR version changes are backward-compatible (new readers can read old
archives); PATCH version changes are bidirectional-compatible.

\textbf{R-13.2 (Forward compatibility).} An implementation
supporting format version $X.Y$ MUST read archives with version
$X.Z$ for any $Z \leq Y$. It MAY refuse to read archives with
version $X.Z$ for $Z > Y$, but if it does read them, MUST honor all
features it understands and reject features it does not.

\textbf{R-13.3 (Predictor compatibility).} A predictor model is
identified by a hash of its model description. Archives include the
predictor hash; readers MUST verify the hash before decoding.
Mismatched hashes MUST cause a hard failure with a clear error
message.

\textbf{R-13.4 (Migration).} Tooling MUST be provided to migrate
archives from version $X.Y$ to $X.Y+1$ if the format change affects
representation. Migration tools MUST be lossless and reversible
where possible.

\textbf{14. Traceability to the main paper}

This section provides the explicit traceability matrix from main-paper
theorems to specification requirements. Reviewers and implementers can
use this to verify that the specification covers every relevant theorem.

\textbf{Theorem 4.1, 4.2 (Type-I achievability and converse)} → Type-I
implementation (separately specified in encoder spec, out of scope
here).

\textbf{Theorem 4.3, 4.4 ($E_{12}$ NP-hard, inapproximable)} →
Heuristic $E_{12}$ selection in encoder; spec does not require
optimal $E_{12}$.

\textbf{Lemma 5.1, Theorems 5.2--5.5, Remark 5.6, Proposition 5.7 (Type-II)} → Type-II implementation
(separately specified in encoder spec).

\textbf{Theorems 6.1, 6.2 (Type-III-A mode selection, tight
$2\varepsilon$ bound)} → Mode-selection logic in encoder
(separately specified).

\textbf{Theorems 6.3, 6.4 (Type-III-B mutual-information and
decomposition invariance)} → Byte tensor coding (encoder spec).

\textbf{Theorems 6.5, 6.6, 6.7, 6.8 (Type-III-C MDL, bits-back, NP-hard,
APX-hard)} → Token tensor coding (encoder spec).

\textbf{Theorem 7.1 (composition correctness)} → §3 (universal
requirements R-3.1 through R-3.8) and §10 (mode-flag stream)
together ensure each mode coder produces a uniquely decodable repair
stream conditional on its mode flag; the composed archive is therefore
uniquely decodable by direct application of the theorem.

\textbf{Theorem 8.1 (no universal dominance)} → Acknowledged in archive
header (fallback-mode flag).

\textbf{Theorem 9.1 (verification)} → Implemented via per-sync-point
hashes (§11.1) or Merkle tree (§11.2).

\textbf{Theorem 10.1 (bit-exact forward pass)} → §3 (universal
requirements R-3.1 through R-3.8), §4 (per-target paths), §5 (integer
arithmetic), §6 (lookup tables).

\textbf{Definition 10.1, Theorem 10.2 (compatible online optimizer)} →
§7 (adaptive predictor conformance).

\textbf{Definition 10.2, 10.3, Theorem 10.3 (random access)} → §8
(random-access conformance with sync-point and seek-index formats).

\textbf{Theorem 10.4 (random access for non-W-bounded predictors)} →
§8.1.2 (state snapshot serialization).

\textbf{Theorem 10.5 (block-device interface)} → §9 (block-device driver
specification).

\textbf{Theorem 10.6 (hardware-target portability)} → §4 (per-target
conformance specifications), §12 (cross-target byte-identity test).

\textbf{15. Closing}

This specification is the executable analogue of the theoretical paper.
Following it produces a reference implementation that satisfies every
theorem proven in the main paper, on every named hardware target. The
conformance test suite of §12 verifies the implementation; the
cross-target byte-identity criterion ensures portability.

The specification is deliberately RFC-style: numbered requirements with
explicit MUST/SHOULD/MAY language, traceability to the underlying
theorems, and a versioning policy. This is the format that allows
independent implementations to verify each other and that supports
long-term archive readability.

Three closing observations. First, every requirement here is traceable
to a theorem in the main paper; no specification choices are arbitrary.
Second, the per-target sections (§4.1 through §4.4) absorb all
hardware-specific complexity, leaving the rest of the specification
target-neutral. Third, the conformance test suite of §12 is the line
between \textquotesingle has been built\textquotesingle{} and
\textquotesingle has been certified\textquotesingle; until an
implementation passes the suite, the deployment claim of §10.8 of the
main paper does not hold for that implementation.

Implementing this specification is engineering work. The mathematical
content is finished. What remains is the disciplined translation of
theorems into code, the verification that the code matches the
specification, and the long-term maintenance of the conformance suite as
hardware targets evolve.

\end{document}
